\section{Related Work}
\label{section:related}

Chapter~\ref{section:background} established the machinery this thesis builds on. This chapter situates \ourloss{} against the specific literature it competes with, organized along the three axes on which it makes a claim: the granularity of the textual reward, the treatment of visual evidence, and the construction of preference data in clinical settings. A fourth section covers the literature that supplies the diagnostic instruments used in the evaluation, and Section~\ref{sec:rel:positioning} collects the comparison into a single table.

\subsection{Fine-grained Alignment in the General Domain}
\label{sec:rel:general}

Traditional DPO suffers from two acknowledged weaknesses: sequence-level sparsity in the reward, and distribution shift between the reference policy and the preference data. A substantial body of work addresses the first.

\textbf{Sub-sequence rewards.} RRPO~\cite{sarkar2025self} restricts the ranking loss to sub-sequence spans, using refined rewards computed only on the tokens that differ between the preferred and rejected responses, together with a forward token-wise KL for stable training. As discussed in Section~\ref{sec:bg:alignment}, the explicit KL term is not an optional addition but a necessary correction: restricting the ranking loss to differing tokens silently removes the implicit KL constraint that DPO inherits through the $\pi_{\text{ref}}$ denominator on all remaining positions. RRPO is the direct antecedent of this work, developed for general-domain video-language alignment, and Chapter~\ref{section:analysis} is in large part an argument about which of its choices survive the move to a clinical setting.

Mask-DPO~\cite{gu2025mask} and ASPO~\cite{wang2025aspo} employ sentence-level mechanisms. Mask-DPO's motivation is factuality: sentences whose correctness cannot be independently verified are masked out so that unverifiable content contributes no gradient, which prevents the objective from being driven by claims the annotation pipeline could not check. This is a reasonable design in the general domain. Section~\ref{sec:res:main} shows it degrades badly on long-form clinical reports, and the ablation in Section~\ref{sec:res:abl-kl} identifies the absent regularizer as the cause rather than the masking itself --- a diagnosis available here only because, in this implementation, MASK-DPO is exactly the $\alpha = 0$ configuration of the same trainer. ASPO instead weights segments adaptively rather than masking them outright, which avoids discarding partial signal but requires a reliable weighting function.

Token-level DPO~\cite{zeng2024token} pushes granularity to its limit by reformulating preference optimization in the token-level Markov decision process, so that credit assignment becomes a per-step quantity derived from the optimization rather than an externally supplied annotation. TGDPO~\cite{zhu2025tgdpo} takes a related route, deriving token-level reward guidance to shape the update.

These approaches differ in where the knowledge of \emph{which tokens matter} comes from. Token-level DPO and TGDPO infer it during optimization; Mask-DPO obtains it from an external factuality check; RRPO derives it from template alignment against a reference. The strategy adopted in Chapter~\ref{section:method} is different again: the identification is made explicit at data construction time, because the preference pair is generated by a minimal edit and the pipeline simply records where the edit occurred. This trades a dependence on inference-time estimation for a dependence on the construction pipeline, which is a real cost --- it is the first limitation acknowledged in Section~\ref{sec:con:limitations} --- but it yields spans that are exact rather than estimated.

\textbf{On-policy construction.} To mitigate the gap between the reference policy and the preference distribution, OPA-DPO~\cite{yang2025opadpo} advocates constructing preferred responses by sampling from the model and refining them, so that the target remains within the model's reachable manifold. Yang et al.~\cite{yang2025mitigating} make the complementary empirical case that on-policy data is the decisive factor in mitigating hallucination through DPO, over and above choices of objective.

This thesis adopts the same principle but reaches it from a different direction, and the difference is not merely presentational. The general-domain argument for on-policy data is about \emph{reachability}: a target the model would essentially never generate cannot be approached by small parameter updates. Section~\ref{sec:ana:reward-hacking} identifies a second and more damaging mechanism specific to the clinical case --- the systematic difference in length and register between clinical reference answers and model outputs constitutes a spurious feature that separates the pair more reliably than the clinical content does. Section~\ref{sec:res:style-gap} measures this directly and finds that $99\%$ of reference-style pairs are separable on response length alone. On-policy construction is therefore not only about staying on the manifold but about removing a confound that would otherwise dominate the learning signal.

What unites this literature is that it focuses primarily on textual signals, neglecting the fine-grained interdependence in which clinical tokens must be strictly tethered to specific visual regions.

\subsection{Multimodal Preference Learning}
\label{sec:rel:multimodal}

Incorporating visual preferences has shown promise in improving grounded interpretation by moving beyond purely textual feedback.

mDPO~\cite{wang2024mdpo} is the foundational instance. It augments the DPO objective with a conditional preference over the same response paired with an original and a perturbed image, so that the objective cannot be satisfied by a policy that ignores the image entirely. This resolves the structural gap noted in Section~\ref{sec:bg:alignment}, where nothing in the standard multimodal DPO objective distinguishes a model that reads the image from one that does not.

CHiP~\cite{fu2025chip} proposes cross-modal hierarchical objectives that optimize textual and visual representations simultaneously and at multiple levels of granularity within each modality, on the argument that alignment is not a single-scale phenomenon. SymPO~\cite{shukla2025sympo} uses symmetric image-text pairs to reinforce the model's reliance on visual evidence over linguistic priors, constructing the preference so that the two modalities must agree.

These approaches share the insight that image-contrast pairs can align the model's multimodal manifold, and \ourloss{} inherits it directly. They differ from this work in two respects, both of which are tested empirically in Chapter~\ref{section:results}.

\textbf{Granularity of the perturbation.} Existing methods typically use coarse image-level perturbations --- cropping, whole-image noise, diffusion corruption, or blacking out the image entirely. These fail to capture the subtle, localized pathological features, such as minute lesions or focal radiographic anomalies, that are decisive in clinical diagnosis. The distinction is not cosmetic: removing the entire image tests whether the model uses vision at all, which is a far weaker property than whether it uses the diagnostically decisive region. A model can retain a globally sensible reading of a chest radiograph while being completely insensitive to the small focal opacity that determines the answer. The ablation in Section~\ref{sec:res:abl-visual} quantifies the gap, showing that replacing lesion-targeted corruption with global noise costs $1.35$ points on average and with cropping $8.05$ points. A second, subtler problem with aggressive global corruption is that it produces an image trivially identifiable as degraded, allowing the model to condition on the corruption itself rather than learning to depend on the lesion.

\textbf{Direction of the preference.} The preference over the perturbed image is generally formulated to prefer the correct response under the clean image over the same response under the corrupted one. This teaches the model that clean evidence is better, but it never teaches it how to behave when evidence is absent. Section~\ref{sec:met:ranking} argues for reversing the preference under corruption instead, so that the objective explicitly prefers the less specific answer when the supporting evidence has been destroyed. Section~\ref{sec:res:abl-visual} tests both alternatives directly and finds the reversed formulation stronger by $3.43$ points over the nearest analogue of the conventional choice.

\subsection{Preference Learning in Clinical Settings}
\label{sec:rel:clinical}

The success of DPO in the general domain has encouraged its adoption in medical workflows, where applications require significantly higher factual precision than general-purpose assistance.

The principal obstacle is the cost of preference labels, and the clinical literature is largely organized around ways of avoiding it. CheXalign~\cite{hein2025chexalign, hein-etal-2025-chexalign} sidesteps human annotation entirely by using automated reference-based scoring of radiology reports to construct preferences, effectively substituting a metric for an annotator. Rad-DPO~\cite{puli2024raddpo} employs specialized data construction to reduce hallucinations in medical VQA. RRG-DPO~\cite{liu2025rrg} targets clinical accuracy in report generation, and R-DPO~\cite{liu2025rdpo} applies recursive optimization to iteratively refine policy factuality, re-deriving preferences from the improved policy. Savage et al.~\cite{savage2025fine} provide a comparative evaluation of supervised fine-tuning against DPO for clinical language models, finding the advantage of preference optimization to be smaller and less consistent than general-domain results would suggest.

MMedPO~\cite{zhu2024mmedpo} is the closest in spirit to this work. It introduces clinical-aware token weighting so that medically salient tokens contribute more to the objective, alongside a multimodal preference component. The difference from the approach taken here is instructive: MMedPO \emph{reweights} a sequence-level reward by estimated clinical salience, whereas \ourloss{} \emph{restricts} the reward to spans identified exactly at construction time. Reweighting retains a contribution from every token and depends on the quality of the salience estimate; restriction discards the uninformative positions entirely but depends on the construction pipeline having identified them correctly.

Two findings from this literature frame the problem this thesis addresses. Kim et al.~\cite{kim2026benchmarking} benchmark DPO across medical LVLMs and report inconsistent gains over supervised fine-tuning --- the empirical anomaly that Chapter~\ref{section:analysis} sets out to explain, and the observation that motivates treating the failure as structural rather than as a matter of insufficient tuning. DeLucia et al.~\cite{delucia2026same} document systematic disagreement between language-model judges and clinicians, which is why substituting AI-generated preference labels is not a satisfactory answer to the annotation cost problem, and why the reference-substitution workaround analyzed in Section~\ref{sec:ana:reward-hacking} became standard in the first place. The same finding is a standing caveat on any evaluation, including this thesis's own, that relies on a model rather than a clinician to judge correctness.

While these methods adapt the data or the weighting, they all rely on the standard DPO schema for reward assignment, which fails to provide the fine-grained supervision required to distinguish between critical clinical findings and generic filler text. None of them addresses the stylistic confound introduced by their own preference construction, and none constrains the policy with an explicit regularizer beyond the implicit one DPO provides.

\subsection{Visual Grounding and Hallucination in Medical LVLMs}
\label{sec:rel:grounding}

A parallel literature characterizes the failures that alignment is meant to fix, and supplies the instruments used to measure them.

MediHall and related work document and categorize hallucination in medical LVLMs~\cite{zhang2024medihall}, establishing that the phenomenon is systematic rather than sporadic. Hallucinogen~\cite{seth2025hallucinogen} benchmarks hallucination arising specifically in implicit reasoning, where the model must combine observations rather than report them. CARES~\cite{xia2024cares} provides a broad trustworthiness benchmark for medical vision-language models spanning factuality, robustness, and instruction adherence, and is a natural target for assessing whether alignment gains generalize beyond the benchmark a model was tuned on.

Methodologically important for this thesis is the work on measuring grounding internally rather than through outputs. Khayatkhoei et al.~\cite{khayatkhoei2025mllms} show that multimodal large language models frequently attend to the correct image region even when they answer incorrectly, and introduce the Attention Ratio as a measure of that attention. The implication is the premise of this thesis's evaluation strategy: output correctness and visual grounding are separable properties, and a method that improves one need not improve the other. Without an internal measurement, an accuracy gain cannot be distinguished from an improvement in linguistic association.

VGMED~\cite{liu2026vgmed} operationalizes this for the medical domain, building over a subset of SLAKE to supply the per-question region annotations the protocol requires and separating clinically critical tokens into localization and attribute categories. Its per-layer formulation is what makes the depth-resolved analysis of Section~\ref{sec:res:grounding} possible.

This literature supplies the diagnosis and the instruments; it does not propose training objectives. \ourloss{} uses its measurements purely for evaluation rather than as a training signal, which is a deliberate choice: it keeps the grounding evaluation independent of the objective being tested, so that the improvement reported in Section~\ref{sec:res:grounding} cannot be an artefact of having optimized the metric.

\subsection{Positioning of This Work}
\label{sec:rel:positioning}

Table~\ref{tab:positioning} summarizes the comparison along the axes developed above. Three of the columns correspond directly to the desiderata of Section~\ref{sec:ana:desiderata}.

\begin{table}[ht]
\centering
\small
\caption{Positioning of \ourloss{} against related preference-optimization methods. \emph{Fine-grained text} indicates reward restricted to sub-sequence spans rather than whole responses. \emph{Visual pref.} indicates an objective term contrasting images. \emph{On-policy} indicates preference pairs derived from the model's own generations. \emph{KL} gives the direction of the explicit regularizer, where one is used: F = forward, R = reverse, F+R = bidirectional, and -- = none beyond DPO's implicit constraint. \emph{Lesion-level} indicates image perturbation targeted at the query-relevant region rather than the whole frame.}
\label{tab:positioning}
\begin{tabular}{llccccc}
\toprule
\textbf{Method} & \textbf{Domain} & \textbf{Fine-grained} & \textbf{Visual} & \textbf{On-} & \textbf{KL} & \textbf{Lesion-} \\
                &                 & \textbf{text}         & \textbf{pref.}  & \textbf{policy} &     & \textbf{level} \\
\midrule
DPO~\cite{rafailov2023direct}        & General  & $\times$   & $\times$   & $\times$   & --  & $\times$ \\
IPO~\cite{azar2024general}           & General  & $\times$   & $\times$   & $\times$   & --  & $\times$ \\
KTO~\cite{ethayarajh2023kto}         & General  & $\times$   & $\times$   & $\times$   & --  & $\times$ \\
Token-DPO~\cite{zeng2024token}       & General  & \cmark     & $\times$   & $\times$   & --  & $\times$ \\
TGDPO~\cite{zhu2025tgdpo}            & General  & \cmark     & $\times$   & $\times$   & --  & $\times$ \\
Mask-DPO~\cite{gu2025mask}           & General  & \cmark     & $\times$   & $\times$   & --  & $\times$ \\
ASPO~\cite{wang2025aspo}             & General  & \cmark     & $\times$   & $\times$   & --  & $\times$ \\
OPA-DPO~\cite{yang2025opadpo}        & General  & $\times$   & $\times$   & \cmark     & --  & $\times$ \\
RRPO~\cite{sarkar2025self}           & Video    & \cmark     & $\times$   & $\times$   & F   & $\times$ \\
\midrule
mDPO~\cite{wang2024mdpo}             & General  & $\times$   & \cmark     & $\times$   & --  & $\times$ \\
CHiP~\cite{fu2025chip}               & General  & \cmark     & \cmark     & $\times$   & --  & $\times$ \\
SymPO~\cite{shukla2025sympo}         & General  & $\times$   & \cmark     & $\times$   & --  & $\times$ \\
\midrule
CheXalign~\cite{hein2025chexalign}   & Clinical & $\times$   & $\times$   & \cmark     & --  & $\times$ \\
Rad-DPO~\cite{puli2024raddpo}        & Clinical & $\times$   & $\times$   & $\times$   & --  & $\times$ \\
RRG-DPO~\cite{liu2025rrg}            & Clinical & $\times$   & $\times$   & $\times$   & --  & $\times$ \\
R-DPO~\cite{liu2025rdpo}             & Clinical & $\times$   & $\times$   & $\times$   & --  & $\times$ \\
MMedPO~\cite{zhu2024mmedpo}          & Clinical & $\times$\textsuperscript{$\dagger$} & \cmark & $\times$ & -- & $\times$ \\
\midrule
\ourloss{} (this work)               & Clinical & \cmark     & \cmark     & \cmark     & F+R & \cmark \\
\bottomrule
\end{tabular}

\vspace{0.4em}
{\footnotesize \textsuperscript{$\dagger$} MMedPO reweights a sequence-level reward by clinical salience rather than restricting the reward to selected spans.}
\end{table}

The table makes the gap explicit. Fine-grained textual reward and multimodal preference each appear in prior work, and on-policy construction appears in a third line of work, but no existing method combines all three. Every clinical method in the table relies on a sequence-level reward, which is the single most consistent finding of this survey and the strongest evidence that the granularity problem has been overlooked rather than solved. No method other than this one uses a bidirectional regularizer --- where an explicit regularizer appears at all it is forward-only --- and none perturbs the image at the granularity of the lesion rather than the frame.

The contribution of this thesis is therefore not any single one of these components but the argument, developed in Chapter~\ref{section:analysis} and tested in Chapter~\ref{section:results}, that in a clinical setting the three failures are coupled rather than independent. A fine-grained reward is undermined by a stylistically confounded preference pair, because the model will exploit the confound before it engages with the spans the reward has been carefully restricted to. A localized reward without a regularizer destabilizes long-form generation, because restricting the ranking loss also removes the implicit constraint that kept the policy near its reference. And neither addresses the possibility that the model is not reading the image at all. Addressing any one in isolation leaves the mechanism intact by another route, which is a plausible explanation for the inconsistent gains that Kim et al.~\cite{kim2026benchmarking} report across the medical DPO literature.
